% ==========================================================
\chapter{代数幾何の記号・用語集}\label{app:H}
% ==========================================================

本編で用いる主な記号と用語をまとめる．「初出」欄は，その記号・用語が最初に定義（導入）される章を示す（「第N章で扱う」といった予告的な言及は除く．本編で定義せず付録で定義するものは付録を挙げる）．
記号の意味は本編の定義が優先する．最後の節では，代数幾何・微分幾何の用語と物理の用語の対応をまとめる．

\newcommand{\appHtablehead}{\textbf{記号} & \textbf{意味} & \textbf{初出}\\ \hline}

% ----------------------------------------------------------
\section{代数的記号}
% ----------------------------------------------------------

集合・写像および一般的な代数構造に関する記号．

\begin{center}
\small
\begin{tabular}{p{0.24\textwidth}p{0.53\textwidth}p{0.12\textwidth}}
\appHtablehead
$\N,\Z,\Q,\R,\C$ & 自然数，整数，有理数，実数，複素数 & 第1章\\
$k$ & 基礎体（多くの場合 $\C$ または代数閉体） & 第1章\\
$\bar k$ & $k$ の代数閉包 & 第5章\\
$\F_p,\ \F_q$ & 位数 $p$（素数），$q$ の有限体 & 第2章\\
$k^\times,\ R^\times$ & 体 $k$ の乗法群，環 $R$ の単元群 & 第2章\\
$\operatorname{char}k$ & $k$ の標数 & 第2章\\
$X/{\sim}$ & 同値関係 $\sim$ による商集合 & 付録A\\
$[x]$ & $x$ の同値類 & 付録A\\
$G/H$，$X/G$ & 剰余類の集合，群作用の軌道空間 & 付録B\\
$\GL_n,\ \SL_n,\ \PGL_n$ & 一般線形群，特殊線形群，射影線形群 & 第1章，第8章\\
$U(n),\ SU(n)$ & ユニタリ群，特殊ユニタリ群 & 付録B\\
$\mu_n$ & $1$ の $n$ 乗根の群 & 第15章\\
$\Hom(V,W)$ & 線形写像（準同型）全体 & 付録C\\
$\End V,\ \Aut V$ & 自己準同型全体，自己同型群 & 付録C\\
$\Ker f,\ \Img f,\ \Coker f$ & 核，像，余核 & 第2章\\
$\rank$ & 行列・線形写像の階数，ベクトル束の階数 & 付録C\\
$\det,\ \Tr$ & 行列式，トレース & 付録C\\
$V^*$ & 双対空間 & 付録C\\
$V\otimes W$ & テンソル積 & 付録C\\
$\mathrm{Sym}^pV,\ \Lambda^pV$ & 対称積，外積 & 付録C\\
$\det V$ & 行列式直線 $\Lambda^{\dim V}V$ & 付録C\\
$[K:k]$ & 体の拡大次数 & 付録B\\
$\mathop{\mathrm{tr.deg}}\nolimits_kK$ & 超越次数 & 第7章\\
$\cong$ & 同型 & 第2章\\
$\simeq$ & 圏同値 & 付録G\\
\end{tabular}
\end{center}

% ----------------------------------------------------------
\section{環論の記号}
% ----------------------------------------------------------

\begin{center}
\small
\begin{tabular}{p{0.24\textwidth}p{0.53\textwidth}p{0.12\textwidth}}
\appHtablehead
$R,\ A$ & 可換環（単位元を持つ） & 第2章\\
$k[x_1,\dots,x_n]$ & $n$ 変数多項式環 & 第1章\\
$k[x_0,\dots,x_n]_d$ & $d$ 次同次多項式の空間 & 第9章\\
$k(x_1,\dots,x_n)$ & 有理関数体 & 第2章\\
$(f_1,\dots,f_r)$ & $f_1,\dots,f_r$ で生成されるイデアル & 第2章\\
$I+J,\ IJ,\ I\cap J$ & イデアルの和，積，共通部分 & 第2章\\
$\af,\ \bfr$ & イデアル & 第2章\\
$\p,\ \q$ & 素イデアル & 第2章\\
$\m$ & 極大イデアル & 第2章\\
$R/I$ & 剰余環（商環） & 第2章\\
$\sqrt I$ & イデアル $I$ の根基 $\{f\mid f^n\in I\}$ & 第2章\\
$\Frac R$ & 整域 $R$ の商体 & 第2章\\
$S^{-1}R$ & 積閉集合 $S$ による局所化 & 第6章\\
$R_f$ & $S=\{1,f,f^2,\dots\}$ による局所化 & 第6章\\
$R_\p$ & 素イデアル $\p$ における局所化 & 第6章\\
$(R,\m)$ & 極大イデアル $\m$ を持つ局所環 & 第6章\\
$\m/\m^2$ & Zariski 余接空間 & 第7章\\
$\dim R$ & Krull 次元 & 第7章\\
$\operatorname{ht}\p$ & 素イデアルの高さ & 第7章\\
$M\otimes_RN$ & $R$-加群のテンソル積 & 付録C\\
\end{tabular}
\end{center}

% ----------------------------------------------------------
\section{幾何学的記号}
% ----------------------------------------------------------

\begin{center}
\small
\begin{tabular}{p{0.24\textwidth}p{0.53\textwidth}p{0.12\textwidth}}
\appHtablehead
$\A^n$ & $n$ 次元アフィン空間 & 第3章\\
$\PP^n$ & $n$ 次元射影空間 & 第8章\\
$[x_0:\dots:x_n]$ & 同次座標 & 第8章\\
$U_i=\{x_i\ne0\}$ & $\PP^n$ の標準的なアフィンチャート & 第8章\\
$\PP(w_0,\dots,w_n)$ & 重み付き射影空間 & 第15章\\
$\Zero(S)$ & 多項式集合 $S$ の零点集合 & 第3章\\
$\Ideal(X)$ & 集合 $X$ の消滅イデアル & 第2章\\
$D(f)$ & 基本開集合 $\{f\ne0\}$ & 第3章\\
$\bar A$ & 閉包（多くの場合 Zariski 閉包） & 第3章\\
$\dim X$ & 次元 & 第7章\\
$\codim_XY$ & 余次元 $\dim X-\dim Y$ & 第7章\\
$T_pX$ & 点 $p$ における（Zariski）接空間 & 第7章\\
$J=(\partial f_i/\partial x_j)$ & Jacobi 行列 & 第7章\\
$\Sing X$ & 特異点集合 & 第7章\\
$\deg X$ & 射影多様体・超曲面の次数 & 第9章\\
$I_p(C,D)$ & 曲線 $C,D$ の点 $p$ における交点重複度 & 第9章\\
$C\cdot D$ & 交点数 & 第9章\\
$\mathrm{Bl}_pX$ & 点 $p$ における blow-up & 第10章\\
$X\times Y,\ X\times_SY$ & 積，ファイバー積 & 第12章\\
\end{tabular}
\end{center}

% ----------------------------------------------------------
\section{多様体の記号}
% ----------------------------------------------------------

代数多様体の射・関数，および複素多様体・ベクトル束に関する記号．

\begin{center}
\small
\begin{tabular}{p{0.24\textwidth}p{0.53\textwidth}p{0.12\textwidth}}
\appHtablehead
$k[X]$ & アフィン多様体 $X$ の座標環 $k[x_1,\dots,x_n]/\Ideal(X)$ & 第3章\\
$\varphi^*$ & 引き戻し（$f\mapsto f\circ\varphi$） & 第5章\\
$k(X)$ & 既約多様体 $X$ の関数体 & 第5章\\
$\OO_{X,p}$ & 点 $p$ における局所環 & 第6章\\
$\m_p$ & $\OO_{X,p}$ の極大イデアル & 第6章\\
$X\dashrightarrow Y$ & 有理写像 & 第10章\\
$\Aut(X)$ & $X$ の自己同型群 & 第10章\\
$z,\ \bar z$ & 複素座標とその複素共役 & 第13章\\
$\partial/\partial z,\ \partial/\partial\bar z$ & Wirtinger 微分 & 第13章\\
$T_pM,\ T_p^*M$ & 接空間，余接空間 & 第7章\\
$TM,\ T^*M$ & 接束，余接束 & 第14章\\
$T^{1,0}X$ & 正則接束 & 第14章\\
$\Omega^p(M)$ & $p$ 次微分形式全体 & 第11章\\
$d,\ \partial,\ \bar\partial$ & 外微分とその $(1,0)$，$(0,1)$ 部分 & 第13章\\
$\pi\colon E\to M$ & ベクトル束（主束） & 第14章\\
$E_p$ & 点 $p$ におけるファイバー & 第14章\\
$g_{\alpha\beta}$ & 変換関数 & 第14章\\
$\Gamma(M,E)$ & $E$ の切断全体 & 第14章\\
$\nabla,\ A$ & 接続（共変微分），接続形式（ゲージ場） & 第14章\\
$F=dA+A\wedge A$ & 曲率（場の強さ） & 第14章\\
$(E,\Phi)$ & Higgs 束 & 第14章\\
$\omega,\ g_{i\bar j}$ & Kähler 形式，Hermite 計量 & 第15章\\
$c_1(E),\ c(E)$ & 第一 Chern 類，全 Chern 類 & 第15章\\
$K_X$ & canonical bundle $\Lambda^nT^{*1,0}X$ & 第14章\\
$\OO(d)$ & $\PP^n$ 上の次数 $d$ の直線束 & 第14章\\
$\Pic X$ & Picard 群（直線束の同型類） & 第11章\\
$\Omega$ & Calabi–Yau 多様体の正則体積形式 & 第15章\\
$h^{p,q}$ & Hodge 数 & 第15章\\
$\chi(X)$ & Euler 標数 & 第13章\\
$\mathcal M$ & モジュライ空間 & 第16章\\
\end{tabular}
\end{center}

% ----------------------------------------------------------
\section{スキームの記号}
% ----------------------------------------------------------

\begin{center}
\small
\begin{tabular}{p{0.24\textwidth}p{0.53\textwidth}p{0.12\textwidth}}
\appHtablehead
$\sheaf F,\ \sheaf G$ & 前層・層 & 第11章\\
$\sheaf F(U)$，$\Gamma(U,\sheaf F)$ & 開集合 $U$ 上の切断 & 第11章\\
$\rho_{UV}$，$s|_V$ & 制限写像 & 第11章\\
$\sheaf F_p$ & 点 $p$ における茎 & 第11章\\
$\OO_X$ & 構造層 & 第11章\\
$\Spec R$ & 環 $R$ の素イデアル全体（素スペクトル） & 第12章\\
$\Zero(I)\subset\Spec R$ & $I$ を含む素イデアル全体（閉集合） & 第12章\\
$D(f)\subset\Spec R$ & $f\notin\p$ なる $\p$ 全体（基本開集合） & 第12章\\
$\kappa(\p)$ & 剰余体 $R_\p/\p R_\p$ & 第6章\\
$\eta$ & 生成点 & 第12章\\
$(X,\OO_X)$ & 環付き空間・スキーム & 第11章，第12章\\
$f^\#$ & （環付き空間・）スキームの射に付随する構造層の射 & 第11章\\
$\Proj S$ & 次数付き環 $S$ の射影スペクトル & 第12章\\
$X_{\mathrm{red}}$ & $X$ の被約化 & 第12章\\
$\Spec k[\varepsilon]/(\varepsilon^2)$ & 双対数の点（接ベクトルを表す） & 第7章\\
$\mathbf{Sch},\ \mathbf{Ring}$ & スキームの圏，可換環の圏 & 付録G\\
$D^b(X)$ & 連接層の有界導来圏 & 第17章\\
\end{tabular}
\end{center}

\subsection*{主な用語}

\begin{description}
\item[代数的集合] 多項式系の共通零点集合．第1章，第3章．
\item[アフィン多様体] 既約なアフィン代数的集合（流儀により代数的集合全体）．第5章．
\item[Zariski 位相] 代数的集合を閉集合とする位相．第3章．
\item[既約] 2つの真の閉部分集合の和に書けないこと．座標環が整域であることと同値．第3章，第5章．
\item[Noether 環] イデアルの昇鎖条件を満たす環．すべてのイデアルが有限生成．第2章．
\item[零点定理] 代数閉体上で $\Ideal(\Zero(I))=\sqrt I$．第4章．
\item[局所化・局所環] 分母を許して点の近傍の情報を取り出す操作と，その結果得られる極大イデアルがただ一つの環．第6章．
\item[正則点・特異点] $\dim T_pX=\dim X$ となる点とそうでない点．第7章．
\item[射・有理写像・双有理] 至る所定義された多項式的写像，稠密開集合上で定義された写像，有理写像による同値．第5章（射），第10章．
\item[前層・層] 開集合ごとのデータと制限写像の組；層はさらに局所性と貼り合わせの公理を満たす．第11章．
\item[スキーム] 局所的にアフィンスキーム $\Spec R$ と同型な局所環付き空間．第12章．
\item[Kähler 多様体] 閉な Kähler 形式を持つ Hermite 多様体．第15章．
\item[Calabi–Yau 多様体] canonical bundle が自明な（コンパクト）Kähler 多様体．Ricci 平坦計量を持つ．第15章．
\item[モジュライ空間] 幾何学的対象の同型類をパラメータ付けする空間．第16章．
\end{description}

% ----------------------------------------------------------
\section{物理との対応用語}
% ----------------------------------------------------------

本編で現れる数学的概念と，物理における対応概念をまとめる．
対応は厳密な同一視ではなく，「物理のこの量はこの数学的構造で記述される」という意味である．

\begin{center}
\small
\begin{tabular}{p{0.36\textwidth}p{0.40\textwidth}p{0.12\textwidth}}
\textbf{数学} & \textbf{物理} & \textbf{初出}\\ \hline
主束 $P\to M$ & ゲージ理論の幾何学的枠組み & 第14章\\
構造群 $G$ & ゲージ群 & 第14章\\
変換関数の取り替え $h_\alpha g_{\alpha\beta}h_\beta^{-1}$ & ゲージ変換 & 第14章\\
主束・ベクトル束の接続 $A$ & ゲージ場（ゲージポテンシャル） & 第14章\\
共変微分 $\nabla=d+A$ & 最小結合をした微分 $D_\mu$ & 第14章\\
曲率 $F=dA+A\wedge A$ & 場の強さ $F_{\mu\nu}$ & 第14章\\
同伴ベクトル束の切断 & 物質場（荷電場） & 第14章\\
Higgs 束の $\Phi\in\Gamma(\End E\otimes K_X)$ & 随伴表現の Higgs 場 & 第14章\\
Chern 類・Chern 数 & インスタントン数，磁荷，位相的量子数 & 第15章\\
平坦接続 & 真空のゲージ配位（ホロノミー） & 第14章\\
変形空間・接空間 $T_p\mathcal M$ & 質量ゼロの揺らぎ（ゼロモード） & 第16章\\
モジュライ空間 $\mathcal M$ & 真空の多様体（真空のモジュライ） & 第16章\\
モジュライ空間の座標 & 質量ゼロスカラー場（モジュライ場）の真空期待値 & 第16章\\
モジュライ空間の次元 & 質量ゼロスカラーの個数 & 第16章\\
モジュライ空間の特異点 & 新たな質量ゼロ粒子の出現，対称性の回復 & 第16章\\
\end{tabular}
\end{center}

\begin{center}
\small
\begin{tabular}{p{0.36\textwidth}p{0.40\textwidth}p{0.12\textwidth}}
\textbf{数学} & \textbf{物理} & \textbf{初出}\\ \hline
Calabi–Yau 3-fold $X$ & 超弦理論の余剰6次元の内部空間 & 第15章\\
Ricci 平坦 Kähler 計量 & 真空の Einstein 方程式の解，超対称性の保存 & 第15章\\
$SU(3)$ ホロノミー & 4次元 $\mathcal N=1$（ヘテロ型）/$\mathcal N=2$（II型）超対称性 & 第15章\\
正則体積形式 $\Omega$ & 共変定数スピノルから作られる形式 & 第15章\\
Kähler モジュライ（$h^{1,1}$ 個） & IIA型でベクトル多重項，IIB型でハイパー多重項 & 第16章\\
複素構造モジュライ（$h^{2,1}$ 個） & IIB型でベクトル多重項，IIA型でハイパー多重項 & 第16章\\
複素化 Kähler 形式 $B+i\omega$ & $B$ 場と体積モジュライの組 & 第16章\\
Euler 標数 $\chi=2(h^{1,1}-h^{2,1})$ & 世代数 $|\chi|/2$（ヘテロ型標準埋め込み） & 第15章\\
Mirror symmetry $h^{1,1}\leftrightarrow h^{2,1}$ & IIA型とIIB型の等価性 & 第17章\\
Blow-up・特異点解消 & 特異点で増大したゲージ対称性の破れ（Coulomb 枝への移行） & 第7章\\
Conifold transition & 質量ゼロブラックホールを介した真空の遷移 & 第15章，第17章\\
部分多様体（＋その上の束） & D-brane（＋その上のゲージ場） & 第17章\\
連接層・導来圏 $D^b(X)$ & B型 D-brane の圏 & 第17章\\
楕円ファイバー空間 & F理論のコンパクト化 & 第17章\\
楕円ファイバーの複素構造 $\tau$ & IIB型の axio-dilaton $C_0+ie^{-\phi}$ & 第17章\\
特異ファイバーの型（Kodaira 分類） & 7-brane 上のゲージ群（ADE 型） & 第17章\\
\end{tabular}
\end{center}

\begin{center}
\small
\begin{tabular}{p{0.36\textwidth}p{0.40\textwidth}p{0.12\textwidth}}
\textbf{数学} & \textbf{物理} & \textbf{初出}\\ \hline
射影空間 $\PP(\mathcal H)$ & 量子状態の空間（純粋状態） & 第8章\\
コンパクト化（射影化） & 無限遠の取り込み，余剰次元の有限化 & 第8章\\
Riemann 面 & 弦の世界面 & 第13章\\
Riemann 面の複素構造のモジュライ & 弦の摂動論的振幅の積分領域 & 第13章\\
正則関数 & カイラル場の量，正則な超ポテンシャル & 第1章\\
代数的集合 $\{\partial W=0\}$ & 超対称真空（F 項条件） & 第1章\\
局所化・局所環 & 点の近傍での局所的な展開（場の理論の局所性とは別） & 第6章\\
層・貼り合わせ & 局所的な場の記述とその大域的整合性 & 第11章\\
冪零元・非被約スキーム & 無限小の揺らぎ，摂動の高次項 & 第12章\\
\end{tabular}
\end{center}

\begin{note}
Kähler モジュライ・複素構造モジュライと超対称多重項の対応は，弦理論の種類（IIA型/IIB型）によって入れ替わる．
これが mirror symmetry の物理的な起源の一つである（第17章）．
また，「局所化」と場の理論の「局所性」は名前が似ているが別の概念であることに注意（第6章）．
\end{note}
